\documentclass[../main.tex]{subfiles}
\begin{document}
%==========================================TIÊU ĐỀ TẬP========================================
\begin{table}[h]
    \begin{adjustwidth}{-1cm}{}
        \begin{tabular}{>{\centering\arraybackslash}p{6cm}|>{\raggedright\arraybackslash}p{12.5cm}}
            \multirow{5}{6cm}
            % Cột bên trái: ÉPISODE và số tập
            {\\[-40pt]
                \flaregothic\fontsize{20pt}{20pt}\selectfont \centering
                \color{eptype}HISTOIRE\color{black}\\
                \freshbot\fontsize{64pt}{64pt}\selectfont
                \color{epnum}A\\[6pt]
                \cafeteria\fontsize{20pt}{20pt}\selectfont\color{epnum}(D-A)\color{black}
                \\[18pt]
                % \sen\fontsize{18pt}{18pt}\selectfont
                % \color{epattr}Interstitiel
                \color{black}
            }
             &
            % Dòng thứ nhất: tên tập tiếng Nhật
            {\fotskip\fontsize{18pt}{18pt}\selectfont \ruby{呼}{よ}び\ruby{名}{な}の\ruby{距}{きょ}\ruby{離}{り}
            }       \\[6pt]
            % Dòng thứ hai: tên tập tiếng Anh
             & {\cafeteria\fontsize{18pt}{18pt}\selectfont A Closer Name}                                  \\[4pt]
            % Dòng thứ ba: tên tập tiếng Việt
             & {\vnmsans\fontsize{18pt}{18pt}\selectfont Khoảng cách của một tiếng gọi}                    \\[8pt]
            % Dòng thứ tư: Nhân vật xuất hiện
             & {\caviar{\fontsize{14pt}{14pt}\selectfont Nhân vật: \emp{kuon1} \emp{ryuuhou1} \emp{aloe}}}
            \\[8pt]
        \end{tabular}
    \end{adjustwidth}
\end{table}
%=========================================TRUYỆN CHÍNH========================================
\color{black}

\txtbx[2]{
    {\color{vol} \uline{Tóm tắt:}} Khi vừa bước sang tuổi mười lăm vào một buổi chiều tháng Tư, Ryuuhou bỗng cảm thấy hai tiếng ``Onii-san'' trang trọng thường ngày dường như không còn phản ánh đúng sự gắn kết giữa hai anh em. Giữa ngôi nhà sáu tầng đang trở thành chốn nương thân ấm áp cho Mano Aloe, nhìn người anh cả Kuon luôn âm thầm gánh vác mọi trách nhiệm từ công ty cho đến gia đình mà chưa từng một lời than vãn, cô bé tomboy quyết định tìm kiếm một cách xưng hô mới. Một tiếng gọi ``Onii'' mộc mạc, ngắn gọn nhưng đong đầy sự thấu hiểu và tin cậy, chính thức xóa nhòa khoảng cách khách sáo để biến nơi đây thành một điểm tựa thực sự của tình thân.
}

\pov{Hikawa Ryuuhou}

Tháng Tư gõ cửa khu phố Anwa bằng những cơn mưa rào đầu mùa dịu mát, gột rửa đi lớp bụi mờ trên những ô cửa kính và mang theo hương đất ẩm thanh lành từ dải công viên ven sông thổi vào.

Vừa bước sang tuổi mười lăm được ít ngày, tôi bỗng thấy bản thân có những thay đổi kỳ lạ mà ngay cả chính mình cũng chẳng thể giải thích nổi. Không phải là chuyện vóc dáng hay sở thích nghe nhạc rock ngổ ngáo; mấy thứ đó thì tôi vẫn thế, vẫn chuộng chiếc áo phông rộng thùng thình, quần lửng túi hộp và mái đầu tém ngắn cũn cỡn chẳng giống thùy mị nết na tí nào.

Cái thay đổi ấy\dots\ nó nằm ở cách tôi nhìn người anh trai của mình.

Chiều thứ Bảy, căn nhà sáu tầng của gia đình chúng tôi chìm trong một sự tĩnh lặng êm đềm. Bố tôi và bác Hijaki sang xưởng cơ khí từ sớm, mẹ thì bận việc dưới kho hàng, còn bà nội đang thong thả ngồi lần tràng hạt ở tầng bốn. Dưới tầng hai, tiếng đàn synthesizer điện tử của Aloe-san thi thoảng lại phát ra vài nốt ngập ngừng rồi tắt hẳn. Cô chị succubus sau bao sóng gió dư luận giờ đây đã tìm thấy một chốn dung thân an toàn, đang cần mẫn nhặt lại từng mảnh vỡ đam mê âm nhạc dưới sự che chở của gia đình tôi.

Tôi ngồi vắt vẻo trên chiếc ghế xoay ở góc phòng khách tầng năm, hai chân co lên đệm, cằm tựa vào đầu gối, đưa mắt nhìn về phía bàn làm việc đối diện.

Ở đó, anh Kuon đang ngồi cặm cụi bên chiếc máy tính xách tay và đống dây cáp mạng ngổn ngang. Chiếc áo thun đỏ giản dị mặc ở nhà đã hơi sờn vai, mái tóc đen bù xù vẫn cứng đầu chĩa vài lọn lên trời. Anh vừa bấm xong đầu nối cáp bằng kìm bấm mạng, vừa đưa tay lên day day thái dương với một tiếng thở dài thật khẽ.

Dưới khóe mắt anh là một quầng thâm nhạt: dấu tích của những đêm thức trắng vừa nghiên cứu tài liệu cho các kỳ thi ở đại học, vừa giải quyết các sự cố kỹ thuật và lịch trình livestream của dàn idol bên văn phòng COVER tại Tokyo.

Anh vừa mới trở về nhà trưa nay sau buổi học ở trường. Vừa đặt chiếc balo nặng trĩu xuống, chưa kịp ngả lưng nghỉ ngơi mười phút, anh đã lập tức xắn tay áo đi kiểm tra lại đường truyền internet ở tầng hai vì sợ cô bé Aloe bị gián đoạn tín hiệu khi tập thu âm.

Lúc nào cũng thế. Lặng lẽ, kiên nhẫn, và chẳng bao giờ biết kêu ca nửa lời.

Tôi mấp máy môi, định cất tiếng hỏi han: ``\textit{Onii-san}\dots''

Nhưng kỳ lạ thay, âm thanh vừa chạm đến đầu lưỡi bỗng khựng lại.

Một cảm giác gợn nhẹ, ngập ngừng dâng lên trong lồng ngực tôi.

``Onii-san''.

Tôi lẩm nhẩm lại cái từ ấy trong đầu, rồi khẽ nhíu mày.

Sao tự dưng nghe nó xa lạ thế nhỉ?

Hồi còn bé tí teo, lúc anh còn sẵn sàng cõng tôi đi dọc con phố ven sông hay chia cho tôi nửa que kem đá, tôi toàn nhõng nhẽo gọi anh là ``Onii-chan''. Cái tiếng gọi trẻ con, ngọt xớt và ấu thơ ấy theo tôi suốt những năm tháng tiểu học. Nhưng rồi khi lớn dần lên, bước vào ngưỡng cửa cấp hai, tôi bắt đầu thấy ngượng. Một đứa con gái tính tình nghịch ngợm, tomboy như tôi mà cứ mở miệng ra là ``Onii-chan'' thì sến súa phát ngượng trước mặt bạn bè. Thế là, như một phản xạ tự nhiên của tuổi mới lớn muốn chứng tỏ mình đã chững chạc, tôi đổi sang gọi anh là ``Onii-san''.

Lịch sự. Đúng mực. Hệt như cách mà người ta vẫn dạy trong các cuốn sách giáo khoa đạo đức. Anh Kuon lúc đầu còn hơi ngỡ ngàng vì thói quen thay đổi đột ngột của tôi, nhưng rồi cũng dần chấp nhận; bản thân anh ấy cũng không có ý kiến về việc này, còn bảo tôi là cứ gọi anh bằng cách nào mà tôi thấy thuận miệng nhất.

Thế nhưng, ngồi nhìn tấm lưng hơi gù xuống vì mệt mỏi của anh lúc này, tôi bỗng nhận ra hai chữ ``Onii-san'' ấy mang một khoảng cách vô hình đến đáng ghét. Cái hậu tố ``-san'' nghe sao mà khách sáo, lễ nghi, hệt như cái cách mà các cô nàng idol ở tận bên Nhật Bản gọi anh mỗi khi có công việc hành chính.

Anh là anh ruột của tôi cơ mà. Là người đã đứng ra che chở cho tôi, là chiếc khiên vững chãi của cả gia đình này, cớ sao tôi lại phải dùng một danh xưng giữ kẽ như thể gọi một người anh họ xa hay một người quen biết ngoài xã hội?

Khoảng cách lễ nghi ấy\dots\ hoàn toàn không đúng với những gì hai anh em tôi đã cùng nhau trải qua.

\ct

``Ryuu này?''

Tiếng gọi trầm ấm bất ngờ kéo tôi giật mình thoát khỏi dòng suy nghĩ mông lung. Tôi chớp mắt lia lịa, bắt gặp ánh mắt đỏ thẫm của anh đang nhìn sang, vẻ mặt lộ rõ nét thắc mắc: ``Em ngồi đực mặt ra đó từ nãy đến giờ làm gì thế? Tivi không bật, điện thoại cũng chẳng buồn lướt. Đói bụng rồi à?''

Tôi vội vàng buông hai chân xuống sàn, hắng giọng một tiếng để lấy lại vẻ tự nhiên: ``Đâu có\dots\ Em đang xem anh bấm dây mạng thôi mà.''

Anh Kuon lắc đầu phì cười, tiện tay quấn gọn cuộn dây cáp màu xám vào cổ tay: ``Bấm xong rồi. Để anh mang xuống tầng hai cắm vào bộ phát phụ cho Aloe-chan. Tiện thể hỏi xem con bé có cần mua thêm gì cho bữa tối không.''

Nhìn anh đứng dậy, với lấy chiếc áo khoác mỏng chuẩn bị bước ra cửa, tôi bất giác đứng bật dậy theo: ``Để em đi cùng! Em mang đĩa táo mẹ gọt sẵn xuống cho cậu ấy luôn.''

Hai anh em cùng bước xuống bậc cầu thang gỗ dẫn về tầng hai. Căn phòng của Aloe nằm ở cuối hành lang, nơi từng là phòng trống dự phòng nay đã được anh Kuon dán thêm các tấm xốp cách âm đơn giản.

Tôi gõ nhẹ ba tiếng lên cánh cửa: ``Aloe ơi, bọn em vào nhé?''

Cánh cửa khẽ mở ra một khe hẹp. Gương mặt nhỏ nhắn với đôi mắt xanh lục bảo của cô nàng Succubus ló ra, trên đầu vẫn đeo chiếc tai nghe kiểm âm lớn. Thấy hai anh em tôi, đôi tai nhọn của chị ấy khẽ động đậy, nở một nụ cười rụt rè nhưng ấm áp: ``Ryuuhou-chan, Kuon-lun\dots\ Mọi người vào đi ạ.''

Căn phòng tầng hai gọn gàng hơn tôi tưởng. Trên bàn làm việc nhỏ là chiếc laptop cũ mà anh Kuon nhượng lại, bên cạnh đặt chiếc micro thu âm có màng lọc âm tự chế.

Anh Kuon nhanh nhẹn quỳ một gối xuống gầm bàn, cắm giắc mạng vào cổng router rồi bật công tắc. Đèn tín hiệu màu xanh lập tức nhấp nháy đều đặn.

``Xong rồi đấy Aloe-chan,'' anh ngẩng lên, mỉm cười xoa nhẹ lớp bụi bám trên đầu ngón tay. ``Đường truyền gigabit nối thẳng từ tầng năm xuống, băng thông ổn định tuyệt đối. Từ giờ em có thể thoải mái tải các tệp beat nhạc nặng mà không sợ bị nghẽn mạng nữa rồi.''

Aloe nhìn những đốm đèn tín hiệu xanh biếc, đôi mắt sáng lên một tia xúc động khó tả. Cô bé bối rối chắp hai tay trước ngực, cúi gập người thật sâu: ``Em\dots\ em cảm ơn anh nhiều lắm ạ! Thật sự phiền anh quá, hôm nay là ngày nghỉ của anh mà\dots''

``Có gì đâu mà phiền,'' anh Kuon xua tay, đứng dậy phủi quần. ``Em cứ an tâm mà luyện thanh với phối nhạc. Ở căn nhà này, không ai phán xét hay làm phiền giấc mơ của em cả. Cứ coi như đây là nhà của mình.''

Tôi đặt đĩa táo giòn ngọt lên bàn cho Aloe, nhìn cô bé cắn một miếng nhỏ với vẻ mặt rạng ngời hạnh phúc. Chờ anh Kuon quay người bước ra hành lang trước, tôi nán lại vài giây, khẽ huých nhẹ vai cô bạn: ``Thế nào? Ở đây quen hơn chưa?''

Aloe ngước nhìn tôi, đôi má hơi ửng hồng. Nàng tiểu quỷ khẽ liếc nhìn bóng lưng cao lớn của anh Kuon đang chậm rãi bước lên cầu thang, rồi thì thầm bằng một chất giọng đầy sự ngưỡng mộ:

``Ryuuhou-chan này\dots\ Anh trai ngươi tuyệt vời thật đấy. Lúc đầu khi mới chuyển tới, ta sợ hãi đến mức chẳng dám bước chân ra khỏi cửa, cứ nghĩ rồi người ta cũng sẽ coi ta như một món hàng lỗi bị vứt bỏ thôi. Nhưng anh ấy chưa từng hỏi một câu nào về quá khứ của ta cả. Cứ lẳng lặng chăm lo, bảo vệ như một lẽ đương nhiên vậy\dots\ Giống như một người anh ruột thực sự.''

Cô bé dừng lại một chút, mắt long lanh nhìn tôi: ``Ngươi thật sự rất may mắn vì có một người anh như thế đấy, Ryuu-chan.''

Câu nói của Aloe như một giọt nước tràn ly, làm rung động toàn bộ những suy nghĩ vốn đang âm ỉ trong lòng tôi suốt cả buổi chiều.

Đúng vậy.

Anh ấy chưa bao giờ xem những gì mình làm là một sự ban ơn. Đối với anh, việc bảo vệ gia đình, chăm lo cho người khác là một điều hiển nhiên đến mức anh sẵn sàng gánh hết mọi nhọc nhằn lên vai mà không màng đòi hỏi hồi đáp.

Vậy thì cớ sao tôi - người em gái ruột thịt được anh yêu thương nhất - lại cứ phải giữ cái khoảng cách khách sáo bằng một tiếng gọi xa xôi?

Tôi bước nhanh lên những bậc cầu thang, trở lại phòng khách tầng năm.

Anh Kuon đang đứng bên bồn rửa nhỏ cạnh gian bếp, cẩn thận xoa bọt xà phòng để rửa sạch những vết dầu mỡ và bụi kim loại bám trên tay lúc sửa máy. Tiếng nước chảy róc rách trong không gian tĩnh mịch. Ánh hoàng hôn vàng óng hắt qua ô cửa sổ kính, dát một mảng sáng ấm áp lên bờ vai anh.

Tôi dừng bước ngay trước ngưỡng cửa bếp, hai tay bất giác siết nhẹ lấy gấu áo hoodie.

Cổ họng tôi khô khốc. Một sự ngập ngừng kỳ lạ lại trỗi dậy, cái cảm giác vừa muốn phá vỡ lớp vỏ bọc cũ, vừa thấy tim đập thình thịch vì ngượng ngùng.

Tôi hít một hơi thật sâu, gom hết tất cả sự can đảm và tính khí thẳng thắn vốn có của mình lại.

Không cần kính ngữ rườm rà. Không cần kiểu cách giữ kẽ. Chỉ là một sự thật trần trụi và thiêng liêng nhất: anh là anh của tôi.

``Này\dots\ \textbf{Onii}.''

Tiếng gọi bật ra khỏi môi tôi. Rất nhỏ. Hơi cộc lốc. Nhưng tròn vành và rõ ràng đến từng âm tiết.

Tiếng vòi nước bỗng chốc ngưng bặt.

Bàn tay đang xoa xà phòng của anh Kuon khựng lại giữa không trung. Lưng anh thẳng lên một nhịp. Căn phòng dường như rơi vào một khoảng lặng kỳ lạ trong chớp mắt, tĩnh lặng đến mức tôi có thể nghe rõ cả tiếng ve sầu sớm râm ran ngoài rặng cây bên kia đường.

Tim tôi đập loạn xạ như đánh trống trận. Chết dở, hay là nghe cộc lốc quá khiến anh giận rồi? Tôi vội vàng siết chặt quai hàm, toan mở miệng chống chế: ``Ý em là\dots''

Thế nhưng, anh Kuon đã từ tốn rút chiếc khăn bông sạch treo bên cạnh, chậm rãi lau khô từng ngón tay. Rồi anh quay đầu lại. Anh không lên tiensg trách móc gì, anh cũng không có nghiêm nghị của một người quản lý thường ngày.

Trên gương mặt hơi mệt mỏi của anh, đôi mắt đỏ thẫm ánh lên sự ngạc nhiên tột cùng, rồi rất nhanh chóng, khóe môi anh từ từ cong lên thành một nụ cười ấm áp và dịu dàng nhất mà tôi từng thấy. Một nụ cười nhẹ nhõm, như thể một tảng băng vô hình vừa tan chảy trong lòng.

``Hửm? Gọi anh có chuyện gì à?'' giọng anh trầm trầm, tự nhiên như hơi thở.

Mặt tôi nóng bừng lên. Biết mình đã vượt qua được cái ranh giới ngượng ngùng ấy, tôi cố gắng giữ vẻ mặt tomboy bất cần đời quen thuộc, khoanh hai tay trước ngực, hất cằm:

``Thì\dots\ em hỏi xem tối nay anh định nấu món gì cho cả nhà với Aloe-san ăn thôi. Đói bụng muốn xỉu rồi đây này.''

Anh Kuon nhìn tôi chăm chú vài giây. Anh không hề gặng hỏi tại sao tôi lại bỏ đi chữ ``-san''. Anh cũng chẳng trêu chọc hay bắt bẻ cái tiếng gọi cộc lốc vừa rồi. Anh hiểu. Giữa hai anh em chúng tôi, có những sự thấu cảm sâu sắc đến mức chẳng cần phải thốt thành lời giải thích.

Anh bật cười thành tiếng, đưa bàn tay ấm áp đã khô ráo lên, vò mạnh mái tóc ngắn của tôi làm nó rối tung lên như một tổ quạ: ``Bố mang cá chép tươi từ chợ về rồi. Tối nay anh nấu canh chua cá, thêm đĩa đậu phụ rán giòn chấm mắm tôm nữa. Em xuống bếp phụ anh nhặt ít rau ngổ với thì là đi.''

``Ui da! Rối hết tóc em rồi!'' tôi phụng phịu gạt tay anh ra, nhưng khóe môi thì không tài nào giấu nổi một nụ cười rạng rỡ. ``Biết rồi, để em nhặt rau cho!''

\ct

Bữa cơm tối hôm ấy trong căn nhà sáu tầng ngập tràn tiếng cười nói rôm rả.

Bên chiếc bàn ăn lớn ở tầng ba, cả gia đình chúng tôi cùng quây quần đông đủ: bố Ryujin, mẹ Kaien, bác Hijaki, bà nội, tôi, anh Kuon và cả cô bé khách trọ nhỏ nhắn Aloe. Nồi canh chua cá nóng hổi bốc khói nghi ngút giữa bàn, mùi thì là, hành hoa và vị chua thanh của mẻ quyện vào nhau thơm lừng cả một góc nhà.

Bố tôi vừa nhấp ngụm bia mát lạnh vừa tấm tắc khen tài nấu nướng của con trai cả, trong khi Aloe thì tròn xoe mắt thích thú thử từng miếng đậu phụ rán giòn rụm chấm nước mắm pha chua ngọt.

Tôi ngồi bên cạnh anh Kuon, thong thả gắp một đũa rau muống luộc, rồi thản nhiên chìa chiếc bát không về phía anh: ``Onii, múc cho em thêm khúc cá nhiều thịt nạc với!''

Bố tôi đang nâng chén bia bỗng khựng lại một giây, mắt liếc nhìn hai anh em đầy thích thú. Mẹ tôi thì mỉm cười hiền từ, còn Aloe chỉ khẽ chớp mắt ngơ ngác nhìn quanh.

Anh Kuon không nói một lời. Anh chỉ mỉm cười, cầm lấy chiếc muôi lớn, khéo léo gạt hết những cọng xương dăm rồi đặt vào bát tôi miếng thịt cá chép béo ngậy, trắng muốt và ngon nhất trong nồi canh.

Tôi cúi đầu và từng miếng cơm nóng hổi, cảm nhận vị ngọt bùi lan tỏa nơi đầu lưỡi. Chỉ là bớt đi một chữ ``san''. Chỉ là một âm tiết ngắn gọn, mộc mạc và có phần cụt ngủn của một đứa con gái tuổi mười lăm.

Thế nhưng, trong khoảnh khắc ấy, tôi biết chắc chắn một điều: đây chính là mái nhà của tôi. Và người con trai đang ngồi lặng lẽ gắp thức ăn cho tôi kia, mãi mãi là điểm tựa vững chãi nhất, là người anh trai duy nhất mà tôi sẵn sàng đứng về phía anh trước mọi giông bão của cuộc đời.

Onii.

Tiếng gọi ấy\dots\ từ nay về sau sẽ là mỏ neo của riêng hai chúng tôi.

\end{document}